% Work-MAT | 2. rocnik | pripravovana hodina 46
% Návaznost: uživatelské připravy.zip (hodiny 01–20), autorská učebnice na GitHubu.
% Realisticky.cz: didaktická inspirace v posloupnosti a typech úloh; zadání jsou vlastní.
\documentclass[11pt,a4paper]{article}
\input{../../sablony/pripravy/priprava.sty}
\nastavPripravu{2. ročník}{Logaritmické rovnice}{46}{Nepřípustné kořeny}
\begin{document}
\blok{Motivace a návaznost}
Na příkladech ukaž, že získaný kořen algebraické rovnice není automaticky kořenem původní logaritmické rovnice. Žáci samostatně porovnají kořeny s podmínkami.
\poznamka{Orientační nabídka pro 45--55 minut: krátké zopakování a motivace, nový výklad, tři společné vzorové úlohy, samostatná práce. Vyber podle tempa třídy; zbývající zadání slouží jako rezerva.}
\blok{Výklad}
Podmínky $F_i(x)>0$ platí pro všechny původní logaritmy zvlášť. Algebraické úpravy nepřípustné kořeny nesmějí vrátit do původního zadání.
\blok{Společné příklady}
\begin{prikladReseniBox}[title={Příklad 1}]
Vyřeš, resp. urči: $\log_2 x+\log_2(x-2)=3$.
\tcblower
\[\begin{aligned}x>2 \\ x^2-2x-8=0\;\Longrightarrow\;x\in\{4,-2\} \\ x=4\end{aligned}\]
\end{prikladReseniBox}
\begin{prikladReseniBox}[title={Příklad 2}]
Vyřeš, resp. urči: $\log_3(x-1)+\log_3(x+1)=1$.
\tcblower
\[\begin{aligned}x>1 \\ x^2=4\;\Longrightarrow\;x\in\{-2,2\} \\ x=2\end{aligned}\]
\end{prikladReseniBox}
\poznamka{Na konci nechej žáky napsat vysvětlení, proč jejich zkouška musí být v původní rovnici.}
\blok{Samostatné procvičování}
\begin{zadaniReseniBox}[title={Samostatná práce 1}]
\begin{enumerate}[label=\alph*)]
\item $\log_{2}x+\log_{2}\left(x-1\right)=\log_{2}\left(6\right)$.
\item $\log_{3}x+\log_{3}\left(x-2\right)=\log_{3}\left(99\right)$.
\item $\log_{5}x+\log_{5}\left(x-3\right)=\log_{5}\left(16000\right)$.
\item $\log_{2}x+\log_{2}\left(x-1\right)=\log_{2}\left(6\right)$.
\item $\log_{3}x+\log_{3}\left(x-2\right)=\log_{3}\left(99\right)$.
\item $\log_{5}x+\log_{5}\left(x-3\right)=\log_{5}\left(16000\right)$.
\end{enumerate}
\tcblower
\begin{enumerate}[label=\alph*)]
\item $x=3;\;\text{druhý algebraický kořen je záporný}$.
\item $x=11;\;\text{druhý algebraický kořen je záporný}$.
\item $x=128;\;\text{druhý algebraický kořen je záporný}$.
\item $x=3;\;\text{druhý algebraický kořen je záporný}$.
\item $x=11;\;\text{druhý algebraický kořen je záporný}$.
\item $x=128;\;\text{druhý algebraický kořen je záporný}$.
\end{enumerate}
\end{zadaniReseniBox}
\blok{Dobrovolné rozšíření}
\begin{zadaniReseniBox}[title={Rozšíření 1}]
\begin{enumerate}[label=\alph*)]
\item $Vytvoř vlastní úlohu tohoto typu a se spolužákem porovnej podmínky a výsledek.$.
\end{enumerate}
\tcblower
\begin{enumerate}[label=\alph*)]
\item Zkontroluj definice a argumenty..
\end{enumerate}
\end{zadaniReseniBox}
\blok{Úkol na doma}
U každého logaritmu vypiš podmínky argumentu a proveď zkoušku. U rovnic vyznač případně nepřípustné kořeny, u nerovnic odpověz intervalem.
\begin{zadaniReseniBox}[title={Domácí úkol}]
\begin{enumerate}[label=\alph*)]
\item $\log_{2}x+\log_{2}\left(x-1\right)=\log_{2}\left(6\right)$.
\item $\log_{3}x+\log_{3}\left(x-2\right)=\log_{3}\left(99\right)$.
\item $\log_{5}x+\log_{5}\left(x-3\right)=\log_{5}\left(16000\right)$.
\item $\log_{2}x+\log_{2}\left(x-1\right)=\log_{2}\left(6\right)$.
\item $\log_{3}x+\log_{3}\left(x-2\right)=\log_{3}\left(99\right)$.
\item $\log_{5}x+\log_{5}\left(x-3\right)=\log_{5}\left(16000\right)$.
\item $\log_{2}x+\log_{2}\left(x-1\right)=\log_{2}\left(6\right)$.
\item $\log_{3}x+\log_{3}\left(x-2\right)=\log_{3}\left(99\right)$.
\end{enumerate}
\tcblower
\begin{enumerate}[label=\alph*)]
\item $x=3;\;\text{druhý algebraický kořen je záporný}$.
\item $x=11;\;\text{druhý algebraický kořen je záporný}$.
\item $x=128;\;\text{druhý algebraický kořen je záporný}$.
\item $x=3;\;\text{druhý algebraický kořen je záporný}$.
\item $x=11;\;\text{druhý algebraický kořen je záporný}$.
\item $x=128;\;\text{druhý algebraický kořen je záporný}$.
\item $x=3;\;\text{druhý algebraický kořen je záporný}$.
\item $x=11;\;\text{druhý algebraický kořen je záporný}$.
\end{enumerate}
\end{zadaniReseniBox}
\poznamka{Domácí úkol obsahuje 8 položek; předpokládaná domácí práce přibližně 15--30 minut. Po větších blocích jsou vědomě promíchané dříve probrané typy.}
\end{document}
